%% LyX 2.4.4 created this file.  For more info, see https://www.lyx.org/.
%% Do not edit unless you really know what you are doing.
\documentclass[english]{article}
\usepackage[T1]{fontenc}
\usepackage[latin9]{inputenc}
\usepackage{enumitem}
\usepackage{amsmath}
\usepackage{amsthm}
\usepackage{geometry}
\geometry{verbose,tmargin=1in,bmargin=1in,lmargin=1in,rmargin=1in}

\makeatletter
%%%%%%%%%%%%%%%%%%%%%%%%%%%%%% Textclass specific LaTeX commands.
\numberwithin{equation}{section}
\numberwithin{figure}{section}
\newlength{\lyxlabelwidth}      % auxiliary length 

%%%%%%%%%%%%%%%%%%%%%%%%%%%%%% User specified LaTeX commands.
\usepackage{fancyhdr}
\usepackage{lastpage}
\pagestyle{fancy}
\usepackage{enumitem}
\usepackage{ifthen}
%sets math fonts to be more readable
\everymath=\expandafter{\the\everymath\displaystyle}
%\DeclareMathSizes{10}{10}{10}{10}
\usepackage{multicol}

\usepackage{tikz}
\usetikzlibrary{plotmarks}
\usepackage{pgfplots}
\pgfplotsset{compat=newest}

\usepackage{calc}
\usepackage[nomessages]{fp}

\usepackage{advdate}    % Advancing/saving dates
\usepackage{datetime}   % Dates formatting
\usepackage{datenumber} % Counters for dates
% Added by lyx2lyx
\setlength{\parskip}{\smallskipamount}
\setlength{\parindent}{0pt}

\makeatother

\usepackage{babel}
\begin{document}
\renewcommand{\author}{Dr.~James Ashe}

\newcommand{\classNum}{137}

\newcommand{\classSec}{B}

\newcommand{\nameType}{Name:}

\newcommand{\examType}{Exam 1}

\renewcommand{\date}{\formatdate{12}{2}{2013}} %%\today

\usdate

%%First page's header%%%%%%%%%%%%%%%%%%%%%%
\fancypagestyle{firststyle} %%header settings for first pagr
{
	\fancyhead[L]{MTH \classNum{}(\classSec{}) \author{}\\
	\textbf{\examType{}} \date{}}
	\fancyhead[C]{\textbf{\nameType}}
	\setlength{\headsep}{14pt}
}
%%%%%%%%%%%%%%%%%%%%%%%%%%%%%%%%%%%%%%%%%%

%%Next page's header%%%%%%%%%%%%%%%%%%%%%%%
\setlist{nolistsep}
\lhead{MTH \classNum{}(\classSec{})} 
\chead{\textbf{\nameType}} 
\cfoot{\thepage{}~of~\pageref{LastPage}} 
\setlength{\headsep}{5pt} 
%%%%%%%%%%%%%%%%%%%%%%%%%%%%%%%%%%%%%%%%%%%

%%Various settings; see the preamble for other settings%%%%%
%\setenumerate[0]{leftmargin=12pt,itemindent=0pt,labelwidth=0pt}
\renewcommand{\labelenumi}{\textbf{\arabic{enumi}.}}
\renewcommand{\labelenumii}{\textbf{\alph{enumii}.}}
\thispagestyle{firststyle}
%%%%%%%%%%%%%%%%%%%%%%%%%%%%%%%%%%%%%%%%%%
\newcommand{\xmax}{0}
\newcommand{\xmin}{-\xmax}
\newcommand{\xstep}{0}
\newcommand{\ymax}{0}
\newcommand{\ymin}{-\ymax}
\newcommand{\ystep}{0} 

\textbf{No credit will be given unless you show your work.} Unless
stated otherwise, each question is worth 5 pts. 

\emph{Solve the following equations.}
\begin{enumerate}[resume]
\item $\frac{7x}{8}-2=3-\frac{15x}{16}$\vfill
\item $4|x+1|-7=6$\vfill
\end{enumerate}
\emph{Use the provide graph to answer the question.}
\begin{enumerate}[resume]
\item The manufacturing cost per vehicle for production of $x$ vehicles
is shown in the graph. Use the graph to estimate the number of vehicles
that must be produced so that the cost per vehicle is \$9,000.\\
\renewcommand{\xmax}{300}
\renewcommand{\xmin}{0}
\renewcommand{\xstep}{0} %must be xmin+n
\renewcommand{\ymax}{40}
\renewcommand{\ymin}{0}
\renewcommand{\ystep}{4} %must be ymin+n
%%%%%%%
\begin{tikzpicture}
	\begin{axis}[
		axis x line=bottom, axis y line=left,     		
		%minor x tick num={0},
		%minor y tick num={0},
		xtick={0,50,...,250},
		ytick={0,10,20,...,40},
		grid=both,,
		%y label style={at={(-0.1,0.65)}},
		xlabel={$x$ Vehicles (thousands)},
		ylabel={Cost each (thousands of \$)},
		xmin=\xmin,xmax=\xmax,
		ymin=\ymin,ymax=\ymax,     		
		width=6cm, height=6cm] 
		\addplot[mark=noner,smooth,red,thick,domain=0:300]{(10^4*x+5*10^8)/(1000000*x)+5}; 
	\end{axis}
\end{tikzpicture}
\end{enumerate}
\emph{Solve the following equation for the provided variable.}
\begin{enumerate}[resume]
\item $G=Q+(M-7)P$ for $M$.\vfill\newpage
\end{enumerate}
\emph{Solve the following word problems}
\begin{enumerate}[resume]
\item Joe bought a pair of high heels for \$322.5 including tax. If the
sales tax is 7.5\% then what was the cost of the high heels before
tax?\vfill
\item Pavlos, the Athenian executioner, needs to make some hemlock poison.
How many gallons of a 80\% hemlock solution must be mixed with 10
gallons of 35\% hemlock to get a mixture that is 65\% hemlock. \vfill
\item John can clean the house in 45 minutes. John can clean the same house
in 15 minutes. How long would it take them to clean the house if they
work together?\vfill\newpage
\end{enumerate}
\emph{Perform the indicated operation(s).}
\begin{enumerate}[resume]
\item Find the center and the radius of the circle. Then graph the circle$(x-2)^{2}+(y+4)^{2}=16$.\\
\renewcommand{\xmax}{10} \renewcommand{\xstep}{5}
\renewcommand{\ymax}{\xmax} \renewcommand{\ystep}{5}
%%%%%%%
\begin{tikzpicture} 
	\begin{axis}[
		axis x line=middle, axis y line=middle,     		
		xtick={-\xmax,-\xstep,...,\xmax},
		ytick={-\ymax,-\ystep,...,\ymax},
		minor x tick num={4},
		minor y tick num={4},
		grid=both,
		xlabel={$$},     
		ylabel={$$},
		xmin=-\xmax.25,xmax=\xmax.25,
		ymin=-\ymax.25,ymax=\ymax.25,     		width=6cm,     		height=6cm] 		%\addplot[mark=noner,smooth,domain=-1:1]{}; 
	\end{axis}
\end{tikzpicture}
\item Write the standard equation for the circle centered at $(9,4)$ and
passing through the origin.\vfill
\item Find the slope, if it exists, of the line passing through the points
$\left(5,0\right)$ and $\left(-2,1\right)$.\vfill
\item Write an equation for the line shown in the graph. \\
\renewcommand{\xmax}{8} \renewcommand{\xstep}{4}
\renewcommand{\ymax}{\xmax} \renewcommand{\ystep}{4}
%%%%%%%
\begin{tikzpicture} 
	\begin{axis}[
		axis x line=middle, axis y line=middle,     		
		xtick={-\xmax,-\xstep,...,\xmax},
		ytick={-\ymax,-\ystep,...,\ymax},
		minor x tick num={3},
		minor y tick num={3},
		grid=both,
		xlabel=$x$,     
		ylabel={$x^2$},
		xmin=-\xmax.25,xmax=\xmax.25,
		ymin=-\ymax.25,ymax=\ymax.25,     		width=6cm,     		height=6cm] 		%\addplot[mark=noner,smooth,domain=-1:1]{}; 
		\addplot[color=red,thick, domain=-\xmax:\xmax] {4/1*(x)-2};
	\end{axis}
\end{tikzpicture}
\item Find an equation for the line which passes through the point $\left(0,1\right)$
and is perpendicular to $y=3x+4$.\vfill\newpage
\end{enumerate}
\emph{Perform the indicated operations and write your answers in the
form $a+bi$.}
\begin{enumerate}[resume]
\item $(2-i)(4+8i)$\vfill
\item $\frac{4}{5+i}$\vfill
\item {[}2 pts{]} $i^{5}$\vspace{1cm}
\item {[}3 pts{]} $i^{24}$\vspace{1cm}
\end{enumerate}
\emph{Use the square root property to find all real or imaginary solutions. }
\begin{enumerate}[resume]
\item $x^{2}-17=0$\vfill
\end{enumerate}
\emph{Solve for $x$.}
\begin{enumerate}[resume]
\item $2x^{2}-36=6x$\vfill
\end{enumerate}
\emph{Solve for $x$. }
\begin{enumerate}[resume]
\item $x^{2}+x-6=0$\\
\renewcommand{\xmax}{8} \renewcommand{\xstep}{4}
\renewcommand{\ymax}{\xmax} \renewcommand{\ystep}{4}
%%%%%%%
\begin{tikzpicture} 
	\begin{axis}[
		axis x line=middle, axis y line=middle,     		
		xtick={-\xmax,-\xstep,...,\xmax},
		ytick={-\ymax,-\ystep,...,\ymax},
		minor x tick num={3},
		minor y tick num={3},
		grid=both,
		xlabel=$x$,     
		ylabel={$x^2+x-6$},
		xmin=-\xmax.25,xmax=\xmax.25,
		ymin=-\ymax.25,ymax=\ymax.25,     		width=6cm,     		height=6cm] 		%\addplot[mark=noner,smooth,domain=-1:1]{}; 
		\addplot[color=red,thick, domain=-\xmax:\xmax] {(x-2)*(x+3)};
	\end{axis}
\end{tikzpicture}
\end{enumerate}

\end{document}
